% Source: https://artofproblemsolving.com/community/c3h1578743_a_silver_issue
% Stable identity: aops-topic-1578743-q01
The first alloy is $\dfrac{2}{5}$ iron and the second is $\dfrac{7}{10}$ iron; the mixture must be $\dfrac{1}{2}$ iron, i.e.\ $15$ pounds of iron in $30$ pounds. With $x$ pounds of the first alloy and $30-x$ of the second,
\[
\dfrac{2}{5}x+\dfrac{7}{10}(30-x)=15\quad\Longrightarrow\quad 4x+210-7x=150\quad\Longrightarrow\quad x=20 .
\]
So use $\boxed{20\text{ pounds of the first alloy and }10\text{ pounds of the second}}$.

Check: iron $8+7=15$ pounds, silver $12+3=15$ pounds.
